\part{Market Failures}

\chapter{Market Failures}
\label{chap:market-failures}

\section{Learning objectives}

After completing this chapter, you should be able to:
\begin{itemize}
  \item explain why a market outcome can fail to maximise social welfare;
  \item analyse externalities using private and social costs and benefits;
  \item classify public goods and common resources; and
  \item identify adverse selection, moral hazard, signalling, and screening.
\end{itemize}

\section{When markets do not deliver efficient outcomes}

Competitive markets can coordinate decentralised decisions effectively, but their efficiency result depends on demanding assumptions. A \keyterm{market failure} arises when private choices do not account properly for all relevant costs, benefits, information, or collective-action problems. This chapter studies three important sources: externalities, public goods and common resources, and information asymmetry.

\section{Externalities}

An \keyterm{externality} occurs when an action affects people not directly involved in a market transaction and the effect is not reflected in the market price. A factory's emissions may harm neighbours; vaccination may reduce infection risk for others. These are respectively negative and positive externalities.

With a negative production externality, marginal social cost exceeds marginal private cost:
\[
MSC=MPC+MEC,
\]
where $MEC$ is marginal external cost. Firms decide using $MPC$, so market output exceeds the efficient level, where marginal social benefit equals marginal social cost. With a positive externality, marginal social benefit exceeds private benefit and market output is too low.

Policy responses include a \keyterm{Pigouvian tax} that internalises external cost, a subsidy for external benefits, direct regulation, and tradable permits. The \keyterm{Coase theorem} shows that clear property rights and costless bargaining can produce efficiency, but transaction costs, many affected parties, power differences, and information problems commonly limit this result.

\section{Public goods and common resources}

A good is \keyterm{rival} if one person's consumption reduces what remains for others, and \keyterm{excludable} if non-payers can be prevented from consuming it. Private goods are rival and excludable. \keyterm{Public goods} are non-rival and non-excludable, such as national defence. \keyterm{Common resources} are rival but non-excludable, such as an open fishery. Club goods are excludable and non-rival up to congestion.

Non-excludability encourages \keyterm{free riding}: people may benefit without contributing, resulting in too little private provision of public goods. Common resources face the opposite problem: individuals have an incentive to use the resource before others do, leading to overuse, the \keyterm{tragedy of the commons}.

Possible responses include public provision, taxation, user fees, property rights, quotas, and community governance. The best institution depends on local knowledge, enforcement capacity, distributional concerns, and the values at stake.

\section{Information asymmetries}

\keyterm{Information asymmetry} exists when one party to a transaction knows something relevant that the other does not. Before a contract, hidden information can create \keyterm{adverse selection}: low-quality sellers or high-risk buyers are disproportionately likely to participate. After a contract, hidden actions can create \keyterm{moral hazard}: insurance may reduce incentives to avoid risk.

A \keyterm{signal} is an action by the informed party that credibly conveys information, such as a warranty. \keyterm{Screening} is an action by the less-informed party to induce self-selection, such as an insurer offering contracts with different deductibles. Reputation, disclosure rules, certification, consumer protection, and contract design can improve outcomes, though they can also impose costs or limit access.

\section{A common lesson}

Market failure does not automatically establish that any government intervention will improve matters. It identifies a divergence between private incentives and a social objective. The next question is comparative institutional analysis: which realistic rule, tax, service, regulation, or private arrangement is most likely to improve the outcome, given information limits and political incentives?

\section*{Chapter Review}

\subsection*{Key concepts}
\begin{itemize}
  \item \textbf{Externality}: an unpriced effect on a third party.
  \item \textbf{Public good}: non-rival and non-excludable good.
  \item \textbf{Common resource}: rival but non-excludable resource.
  \item \textbf{Free riding}: benefiting without contributing.
  \item \textbf{Adverse selection}: hidden information before contracting.
  \item \textbf{Moral hazard}: hidden action after contracting.
\end{itemize}

\section*{Exercises}

\subsection*{Conceptual questions}
\begin{enumerate}[label=\arabic*.]
  \item Define a market failure. Why does its existence not automatically imply that intervention will improve outcomes?
  \item Distinguish a negative externality from a public good.
  \item Explain the difference between a common resource and a public good.
  \item Distinguish adverse selection from moral hazard using insurance examples.
\end{enumerate}

\subsection*{Graphical exercises}
\begin{enumerate}[resume,label=\arabic*.]
  \item Draw demand, marginal private cost, and marginal social cost for a negative production externality. Mark market output, efficient output, and deadweight loss.
  \item Draw a simple diagram showing why private willingness to pay can understate the social benefit of a public good.
  \item Draw a used-car market and explain how uncertainty about quality may reduce trade.
\end{enumerate}

\subsection*{Analytical exercises}
\begin{enumerate}[resume,label=\arabic*.]
  \item Demand is $P=100-Q$, private marginal cost is $20$, and external marginal cost is $20$. Find market and efficient quantities and the corrective tax.
  \item Explain how transferable fishing quotas can reduce overfishing. What distributional issue can arise when quotas are initially allocated?
  \item Explain how a warranty can reduce adverse selection in a used-goods market.
  \item Compare a carbon tax, direct regulation, and no policy in terms of information requirements and incentives.
\end{enumerate}

\subsection*{Exam-style multiple-choice questions}
\begin{enumerate}[resume,label=\arabic*.]
  \item A negative externality exists when: (a) private cost exceeds social cost; (b) social cost exceeds private cost; (c) all firms make losses; or (d) price is zero.
  \item A public good is: (a) rival and excludable; (b) non-rival and non-excludable; (c) rival and non-excludable; or (d) supplied by government by definition.
  \item Hidden information before a contract is associated with: (a) moral hazard; (b) adverse selection; (c) perfect competition; or (d) price-taking.
\end{enumerate}

\subsection*{PPE discussion question}

\noindent Choose one current policy problem---climate change, health insurance, housing, or fisheries. Identify the market failure and compare two realistic institutional responses in terms of efficiency, equity, and political feasibility.
